% !TEX root = ../report.tex
\chapter*{РЕФЕРАТ}
Отчёт содержит \pageref{LastPage} с., \totalfigures{} рис., \totaltables{} табл., 2 источника.

ПОИСК В МАССИВЕ, ЛИНЕЙНЫЙ ПОИСК, БИНАРНЫЙ ПОИСК, РЕКУРСИЯ, ЧИСЛО СРАВНЕНИЙ, ДОПОЛНИТЕЛЬНАЯ ПАМЯТЬ

Объект исследования --- алгоритмы поиска целого числа в массиве. Цель работы --- сравнительный анализ ресурсной эффективности линейного, двух нерекурсивных бинарных и рекурсивного бинарного алгоритмов поиска.

Разработаны схемы четырёх алгоритмов, реализована программа на языке C, выполнено функциональное тестирование. Временная эффективность оценена по числу сравнений искомого значения с элементами массива; дополнительная память оценена теоретически. Экспериментальные минимумы и максимумы числа сравнений получены поиском значений, которые присутствуют и отсутствуют в массиве.

Линейный поиск выполняет от одного до $N$ сравнений и не требует сортировки. В худшем случае число сравнений бинарных алгоритмов растёт логарифмически. Эти алгоритмы требуют упорядоченного массива. Рекурсивный вариант использует дополнительную память для стека вызовов. Алгоритм без ветвления уменьшает шаг поиска вдвое; число его сравнений зависит от размера массива и искомого значения.
